\documentclass[12pt]{article}
\usepackage{makecell}
\usepackage{hyphenat}
\usepackage[utf8]{vietnam}
\usepackage[vietnamese=nohyphenation]{hyphsubst}
\usepackage[vietnamese]{babel}
\usepackage[table,dvipsnames]{xcolor} % Required for row coloring
\usepackage{booktabs} % Required for nicer horizontal lines
\usepackage{ragged2e}
\usepackage{caption}
\usepackage{subcaption}
\usepackage{float}
\usepackage{amsmath,amssymb,amsfonts}
\usepackage{graphicx}
\usepackage{adjustbox}
\usepackage{multicol}
\usepackage[inline,shortlabels]{enumitem}
\usepackage{setspace}
\usepackage{longtable}
\usepackage{bm} 
\usepackage[inkscapelatex=false]{svg}
\usepackage{titling}
\usepackage{tabularx}
\usepackage[T1]{fontenc}
\usepackage{lmodern,mathrsfs}
% \usepackage[draft]{hyperref}
\usepackage{hyperref}
\usepackage{xparse}
\usepackage{fontawesome5}
\usepackage[most]{tcolorbox}
\usepackage{tikz}
\usepackage[framemethod=TikZ]{mdframed}
\usetikzlibrary{matrix, fit, shadows, positioning, arrows.meta, calc, shapes.geometric, backgrounds, shapes.symbols}
\usepackage{paracol}
\usepackage{bbm}
\usepackage{pgfplots}
\usepackage{siunitx}
\usepackage{multirow}
%  biblatex (thay cho IEEEtran.bst) ------------------------------------
\usepackage[
    backend=bibtex,
    style=ieee,
    citestyle=numeric-comp,
    maxnames=1,     % chỉ hiển thị 1 tên tác giả, phần còn lại thành ``et al."
    minnames=1,
    uniquelist=false,
    hyperref=false
]{biblatex}
\usepackage{forest}
\addbibresource{references.bib}
\DefineBibliographyStrings{english}{
  andothers = {et\addabbrvspace al\adddot}
}
\AtBeginBibliography{\selectlanguage{english}}

% -----------------------------------------------------------------
\AtBeginDocument{\renewcommand\refname{}}

\usepackage[toc,page]{appendix}
\renewcommand{\appendixtocname}{Phụ lục}      
\renewcommand{\appendixpagename}{Phụ lục}

\usepackage{xurl}
\setlength{\headheight}{14.5pt}
\setlength{\topmargin}{-.5in} \setlength{\textheight}{9.25in}
\setlength{\oddsidemargin}{-0.1in} \setlength{\textwidth}{6.8in}

\usepackage{tocloft}
\renewcommand{\cfttoctitlefont}{\Huge\bfseries}
\setlength{\cftaftertoctitleskip}{2em}
\setlength{\cftbeforesubsecskip}{7pt} % giãn cách subitem mục lục
\renewcommand{\cftsecleader}{\cftdotfill{\cftdotsep}}
\setlength{\cftsecnumwidth}{3em}
\renewcommand{\cftsecaftersnum}{\quad}
\renewcommand{\cftsubsecleader}{\cftdotfill{\cftdotsep}}
\setlength{\cftsubsecnumwidth}{4em}
\renewcommand{\cftsubsecaftersnum}{\quad}
\renewcommand{\cftsubsubsecleader}{\cftdotfill{\cftdotsep}}
\setlength   {\cftsubsubsecnumwidth}{5em}  
\renewcommand{\cftsubsubsecaftersnum}{\quad}

\PassOptionsToPackage{vietnamese.licr}{babel}
\usepackage{vipythonhighlight}
\usepackage{tvietlistings}
\newcounter{mycounter}
\newcommand\showmycounter{\stepcounter{mycounter}\themycounter}
\usepackage{lipsum}
\newcommand\showlips{\stepcounter{mycounter}\lipsum[\value{mycounter}]}

%--------------------------------------------------------------------
\usepackage{framed}
\usepackage{fancyhdr}
\usepackage{upquote}
\usepackage{listings}	 
	    

\usepackage[toc, nonumberlist, style=super]{glossaries}
\makenoidxglossaries

\newglossarystyle{clean3col}{%
    % 1) Table layout: 3 columns
    \renewenvironment{theglossary}%
    {%
      \begin{longtable}{@{}
        >{\setlength{\parindent}{0pt}\RaggedRight\arraybackslash}p{0.2\linewidth}
        >{\setlength{\parindent}{0pt}\RaggedRight\arraybackslash}p{0.23\linewidth}
        >{\setlength{\parindent}{0pt}\RaggedRight\arraybackslash}p{0.52\linewidth}
      @{}}%
    }%
    {\end{longtable}}%

    % 2) Header row
    \renewcommand{\glossaryheader}{%
        \bfseries Thuật ngữ &
        \bfseries Tạm dịch &
        \bfseries Mô tả \\
        \midrule
        \endhead
    }%

    % 3) Each entry row
    \renewcommand{\glossentry}[2]{%
        \itshape \glstarget{##1}{\glossentryname{##1}} &
        \glsentryuseri{##1} &
        \glossentrydesc{##1} \\
        \addlinespace[6pt]
    }%
}


% Color --------------------------------------------------------------------
\usepackage{epigraph}
\usepackage{titlesec}

\definecolor{codebackground}{RGB}{250, 250, 250}
\definecolor{codecomment}{RGB}{34, 139, 34}
\definecolor{codekeyword}{RGB}{0, 0, 255}
\definecolor{codestring}{RGB}{255, 0, 0}
\definecolor{codeidentifier}{RGB}{0, 0, 0}
\definecolor{codelinenumbers}{RGB}{90, 90, 90}

\definecolor{codegreen}{rgb}{0,0.6,0}
\definecolor{codegray}{rgb}{0.5,0.5,0.5}
\definecolor{codepurple}{rgb}{0.58,0,0.82}
\definecolor{backcolour}{rgb}{0.95,0.95,0.92}

\definecolor{aivnGreen}{RGB}{52,79,30}  
\definecolor{aivnCodeOutputColor}{RGB}{11,83,148}  
\definecolor{captiongray}{gray}{0.15} 
\definecolor{captionbg}{RGB}{250,250,250}

\definecolor{foldercolor}{RGB}{0, 102, 204}
\definecolor{filecolor}{RGB}{104, 104, 104}
\definecolor{scriptcolor}{RGB}{180, 50, 50}
\definecolor{pythoncolor}{RGB}{0, 128, 128}

% --- Extra Brand colors for the table ---
\definecolor{brandIndigo}{HTML}{333D91} 
\definecolor{rowLavender}{HTML}{EEECFA}
\setlength{\parskip}{0.3em}

% Table
\newtcolorbox{captionbelowboxaivn}{
    colback=aivnGreen,
    colframe=aivnGreen,
    boxrule=0pt,
    arc=4pt,
    left=4pt,
    right=4pt,
    top=2pt,
    bottom=2pt,
    width=\linewidth,
    boxsep=4pt,
    halign=center,
    % fontupper=\bfseries\color{white}
    fontupper=\color{white}
}

% Figure
\newtcolorbox{figurecardcaptionaivn}{
  colback=aivnGreen,        % Màu nền (background color)
  colframe=aivnGreen,       % Màu viền box (border color)
  boxrule=0pt,              % Độ dày viền box (0pt = không viền)
  arc=2pt,                  % Độ bo góc (2pt là bo nhẹ, tăng để bo tròn hơn)
  top=3pt,                  % Padding phía trên nội dung
  bottom=3pt,               % Padding phía dưới nội dung
  left=10pt,                % Padding bên trái
  right=10pt,               % Padding bên phải
  width=\linewidth,         % Chiều rộng bằng toàn bộ dòng
  fontupper=\small\fontsize{11pt}{13pt}\selectfont\color{white} 
}

% Tạo counter riêng
\newcounter{hinh}
% Định nghĩa macro caption Hình
\newcommand{\capfig}[1]{%
  \refstepcounter{hinh}%
  % \textbf{Hình~\thehinh:}~#1
    Hình~\thehinh:~#1
}

\newtcolorbox{figurecaptionmodern}{
  colback=captionbg,
  colframe=aivnGreen,
  boxrule=1.5pt,
  bottomrule=2pt,  
  toptitle=0pt,
  bottomtitle=0pt,
  left=6pt,
  right=6pt,
  top=4pt,
  bottom=2pt,
  width=\linewidth,
  fontupper=\small\fontsize{11pt}{13pt}\selectfont\color{black},
  halign=center
}

% Code Block -----------------------------------------------------------------
\newtcolorbox{aivncodebox}[1][]{%
  enhanced,
  colback=gray!10,
  colframe=aivnGreen,
  coltitle=white,
  colbacktitle=aivnGreen,
  title=#1,
  fonttitle=\bfseries\sffamily,
  listing only,
  breakable,
  listing engine=listings,
  listing options={
    language=Python,
    basicstyle=\ttfamily\footnotesize,
    commentstyle=\color{ForestGreen},
    keywordstyle=\color{blue},
    stringstyle=\color{red!70!black},
    numbers=left, upquote=true,
    numberstyle=\tiny\color{gray},
    numbersep=5pt,
    breaklines=true,
    showstringspaces=false,
    xleftmargin=1.5em,
    framexleftmargin=1.5em
  }
}

% Code Output Block -----------------------------------------------------------------
\newtcolorbox{aivnCodeOutput}[1][]{%
  enhanced,
  colback=gray!10,
  colframe=aivnCodeOutputColor,
  coltitle=white,
  colbacktitle=aivnCodeOutputColor,
  title=#1,
  fonttitle=\bfseries\sffamily,
  listing only,
  breakable,
  listing engine=listings,
  listing options={
    language=Python,
    basicstyle=\ttfamily\footnotesize,
    commentstyle=\color{ForestGreen},
    keywordstyle=\color{blue},
    stringstyle=\color{red!70!black},
    numbers=left,
    numberstyle=\tiny\color{gray},
    numbersep=5pt,
    breaklines=true,
    showstringspaces=false,
    xleftmargin=1.5em,
    framexleftmargin=1.5em
  }
}

% Last argument allows more tcolorbox options to be added

% --- Guideline/Tip box ---
\newtcolorbox{guidelinebox}[2][]{
    enhanced, colback=yellow!10!white, colframe=orange!90!black, boxrule=1pt, arc=4pt,
    fonttitle=\bfseries, attach boxed title to top left={yshift=-2mm, xshift=3mm},
    colbacktitle=orange!90!black, coltitle=white,
    title=\faLightbulb\hspace{1mm} #2, #1
}

\newtcolorbox{notebox}[1][]{
  enhanced,  colback=yellow!10, colframe=orange!80, boxrule=1pt, 
  arc=4pt, fontupper=\small, attach boxed title to top left={yshift=-2mm, xshift=3mm},
  colbacktitle=orange!80,  coltitle=white, fonttitle=\bfseries,  
  title=\faInfoCircle~\textbf{Lưu ý},  #1
}

\newtcolorbox{infobox}[1]{%
    enhanced, 
    breakable, 
    colback=green!3, 
    colframe=green!3, 
    arc=5pt,
    borderline west={3pt}{0pt}{aivnGreen},
    attach boxed title to top left={ 
        yshift*=-\tcboxedtitleheight/2, 
        xshift=-0.1em 
    },
    colbacktitle=aivnGreen, %
    coltitle=white, 
    fonttitle=\bfseries, 
    boxed title style={ 
        arc=4pt, 
        top=2pt, 
        bottom=2pt,
        left=8pt,
        right=8pt,
    },
    title={\faInfoCircle\quad #1}
}
\newtcolorbox{examplebox}[1]{%
    enhanced,
    breakable,
    colback=brandIndigo!5,      
    colframe=brandIndigo!5,     
    arc=5pt,
    borderline west={3pt}{0pt}{brandIndigo}, 
    attach boxed title to top left={
        yshift*=-\tcboxedtitleheight/2,
        xshift=-0.1em
    },
    colbacktitle=brandIndigo,   
    coltitle=white,
    fonttitle=\bfseries,
    boxed title style={
        arc=4pt,
        top=2pt,
        bottom=2pt,
        left=8pt,
        right=8pt,
    },
    title={\faLightbulb\quad #1} 
}
% Prompt Block ---------------------------------------------------------------
\definecolor{promptBlue}{HTML}{2F5D8C}
\definecolor{promptBg}{HTML}{F6F9FC}
\definecolor{promptBorder}{HTML}{B7C9DC}

% Dùng cho prompt đã viết theo cú pháp LaTeX bình thường
\newtcolorbox{promptbox}[1][Prompt]{%
    enhanced,
    breakable,
    colback=promptBg,
    colframe=promptBorder,
    coltitle=white,
    colbacktitle=promptBlue,
    fonttitle=\bfseries\sffamily,
    title={#1},
    boxrule=0.8pt,
    arc=4pt,
    left=8pt,
    right=8pt,
    top=2pt,
    bottom=1pt,
    before skip=8pt,
    after skip=8pt,
    borderline west={3pt}{0pt}{promptBlue},
    fontupper=\small\sffamily,
}

% Dùng cho prompt raw: giữ nguyên xuống dòng, không highlight keyword
\newtcblisting{promptboxraw}[1][Prompt]{%
    enhanced,
    breakable,
    listing only,
    listing engine=listings,
    colback=promptBg,
    colframe=promptBorder,
    coltitle=white,
    colbacktitle=promptBlue,
    fonttitle=\bfseries\sffamily,
    title={#1},
    boxrule=0.8pt,
    arc=4pt,
    left=8pt,
    right=8pt,
    top=-4pt,
    bottom=-4pt,
    before skip=4pt,
    after skip=4pt,
    borderline west={3pt}{0pt}{promptBlue},
    listing options={
        language={},
        basicstyle=\small\ttfamily,
        columns=fullflexible,
        keepspaces=true,
        breaklines=true,
        showstringspaces=false,
        numbers=none,
        frame=none,
        keywordstyle=\color{black},
        commentstyle=\color{black},
        stringstyle=\color{black},
        identifierstyle=\color{black}
    }
}


%--------------------------------------------------------------------
\lstdefinestyle{mystyle}{
    backgroundcolor=\color{codebackground},
    commentstyle=\color{codecomment},
    keywordstyle=\color{codekeyword},
    stringstyle=\color{codestring},
    identifierstyle=\color{codeidentifier},
    basicstyle=\ttfamily\footnotesize,
    breakatwhitespace=false,
    breaklines=true,
    captionpos=b,
    keepspaces=true,
    numbers=left,
    numberstyle=\tiny\color{codelinenumbers}, 
    numbersep=5pt,
    showspaces=false,
    showstringspaces=false,
    showtabs=false,
    tabsize=2
}

\lstset{style=mystyle}

\hypersetup{
    colorlinks=true,
    linkcolor=black,
    filecolor=magenta,    
    citecolor=green,
    urlcolor=blue,
    pdftitle={Quiz: Data Warehouse, Data Lake và Data Lakehouse},
    pdfpagemode=FullScreen,
}

%--------------------------------------------------------------------
\newcounter{choice}
\renewcommand\thechoice{($\alph{choice}$)}
\newcommand\choicelabel{\thechoice}

%--------------------------------------------------------------------
\newenvironment{choices}%
  {\list{\choicelabel}%
     {\usecounter{choice}\def\makelabel##1{\hss\llap{##1}}%
       \settowidth{\leftmargin}{W.\hskip\labelsep\hskip 0 em}%
       \def\choice{%
         \item
       } % choice
       \labelwidth\leftmargin\advance\labelwidth-\labelsep
       \topsep=0pt
       \partopsep=0pt
     }%
  }%
  {\endlist}
  
%--------------------------------------------------------------------
\def\CChoice{%
      \choice
      \addanswer{\theenumi}{\thechoice}%
    }
    \let\CChoice\CChoice
%    \par % Uncomment this to have choices always start a new line
   % \let\par\@empty
    % If we're continuing the textbf containing the question,
    % then leave a bit of space before the first choice:
    \ifvmode\else\enskip\fi
    \ignorespaces
  %
  {}
\makeatother
\newbox\allanswers
\setbox\allanswers=\hbox{}
\newcommand{\addanswer}[2]{%
  \global\setbox\allanswers=\hbox{\unhbox\allanswers \quad #1.~#2}%
}
\newcommand{\showanswers}{%
  \vfill
  \begin{center}
    Answers
  \end{center}
  \unhbox\allanswers
}

\makeatletter
\newenvironment{multichoices}[1][1]{%
\begin{multicols}{#1}}{%
\end{multicols}}
\makeatother

\pagestyle{fancy}

\DeclareMathOperator*{\argmax}{arg\,max}
\DeclareMathOperator*{\argmin}{arg\,min}
\DeclareMathOperator{\E}{\mathbb{E}}

% Outline --------------------------------------------------------------------
% Section numbering ----------------------------------------------------------
\renewcommand{\thesection}{\Roman{section}.} 
\renewcommand{\thesubsection}{\thesection\arabic{subsection}.}
\renewcommand{\thesubsubsection}{\thesubsection\arabic{subsubsection}.}

% Paper-like section title sizes --------------------------------------------
% Title paper: 17pt
% Section should be clearly smaller than title but still prominent.
\titleformat{\section}
  {\normalfont\fontsize{16pt}{20pt}\selectfont\bfseries}
  {\thesection}
  {0.85em}
  {}

\titleformat{\subsection}
  {\normalfont\fontsize{14pt}{16pt}\selectfont\bfseries}
  {\thesubsection}
  {0.75em}
  {}

\titleformat{\subsubsection}
  {\normalfont\fontsize{12pt}{15pt}\selectfont\bfseries}
  {\thesubsubsection}
  {0.65em}
  {}

% Paper-like compact spacing -------------------------------------------------
\titlespacing*{\section}
  {0pt}
  {1.25em plus 0.2em minus 0.1em}
  {0.65em plus 0.1em minus 0.1em}

\titlespacing*{\subsection}
  {0pt}
  {1.0em plus 0.2em minus 0.1em}
  {0.45em plus 0.1em minus 0.1em}

\titlespacing*{\subsubsection}
  {0pt}
  {0.85em plus 0.15em minus 0.1em}
  {0.35em plus 0.1em minus 0.1em}


% Logo + Title ---------------------------------------------------------------

\setlength{\droptitle}{-7em}

% Title horizontal rule
\newcommand{\TitleRule}{%
  \rule{0.90\textwidth}{0.6pt}
}

% Chiều cao vùng title giữa 2 rule
\newlength{\TitleBlockHeight}
\setlength{\TitleBlockHeight}{2.2cm}

\pretitle{%
  

  % Logo row
  \begin{tabular}{@{}c@{\hspace{1.0em}}c@{\hspace{1.0em}}c@{}}
    \raisebox{-0.35\height}{\includegraphics[height=2.0cm]{Figures/aio2026_logo.pdf}}    
  \end{tabular}

  %\vspace{0.2em}

\begin{center}
  \TitleRule
  \par\nointerlineskip
}

\title{%
  \parbox[c][\TitleBlockHeight][c]{0.90\textwidth}{%
    \centering
    {\fontsize{20pt}{24pt}\selectfont\bfseries Quiz: Data Warehouse, Data Lake\\ và Data Lakehouse}\\[6pt]
    {\fontsize{13pt}{15pt}\selectfont (Data Architecture)}
  }%
}

\posttitle{%
  \par\nointerlineskip
  \TitleRule
  \end{center}
}

\preauthor{%
  \begin{center}
  \fontsize{11pt}{13pt}\selectfont
}

\author{%
  \textbf{Tác giả 1}
  \hspace{2.8em}
  \textbf{Tác giả 2}\\
  \vspace{1em}
  \textbf{Proofreading: }
}

\postauthor{%
  \par
  \end{center}
}

\date{}
% Header --------------------------------------------------------------------
\pagestyle{fancy}
\fancyhf{}
\lhead{\href{https://www.facebook.com/aivietnam.edu.vn}{\textcolor{aivnGreen}{\bfseries AI VIET NAM (AIO2026)}}}
\rhead{\href{https://aivietnam.edu.vn/}{\textcolor{aivnGreen}{\bfseries aivietnam.edu.vn}}}
\fancyfoot[r]{\textcolor{aivnGreen}{\bfseries\foreignlanguage{vietnamese}{Trang \thepage}}}
\fancyfoot[l]{%
  \href{https://zalo.me/84911118758}{%
    \textcolor{aivnGreen}{\bfseries\foreignlanguage{vietnamese}{SĐT/Zalo: 0911118758 (Dr. Hồng Phúc)}}%
  }%
}
\renewcommand{\headrulewidth}{1.0pt}
\renewcommand{\headrule}{{\color{aivnGreen}\hrule width\headwidth height\headrulewidth \vskip-\headrulewidth}}

\renewcommand{\footrulewidth}{1.0pt}
\renewcommand{\footrule}{{\color{aivnGreen}\hrule width\headwidth height\headrulewidth \vskip-\headrulewidth}}


\fancypagestyle{firstpage}{
    \fancyhf{}
    \fancyfoot[r]{\textcolor{aivnGreen}{\bfseries\foreignlanguage{vietnamese}{Trang \thepage}}}
    \fancyfoot[l]{%
      \href{https://zalo.me/84911118758}{%
        \textcolor{aivnGreen}{\bfseries\foreignlanguage{vietnamese}{SĐT/Zalo: 0911118758 (Dr. Hồng Phúc)}}%
      }%
    }
    \renewcommand{\headrulewidth}{0pt}
    \renewcommand{\footrulewidth}{1.0pt}
    \renewcommand{\footrule}{%
        {\color{aivnGreen}\hrule width\headwidth height\footrulewidth}
    }
}



\begin{document}
\maketitle
\thispagestyle{firstpage}
\vspace{-3em}
\noindent Mỗi câu hỏi có đúng một đáp án đúng. Các số liệu trong đề được lấy từ phần thực hành của bài đọc.

\begin{enumerate}[label=\textbf{\arabic*.}, leftmargin=1.6em, itemsep=1.2em]
\item Một công ty lưu dữ liệu đơn hàng thành các file Parquet dưới prefix \texttt{s3://shop/orders/}. Spark và Trino đọc ``bảng'' \texttt{orders} bằng cách liệt kê mọi file dưới prefix này rồi tự suy ra schema. Hệ thống không có catalog hay transaction log nào. Theo tiêu chí ``lớp nào sở hữu khái niệm bảng'', đây là kiến trúc gì?

\begin{enumerate}[label=\alph*), leftmargin=2.2em, itemsep=1pt]
    \item Data Warehouse, vì dữ liệu được lưu ở định dạng columnar.
    \item Data Lakehouse, vì nhiều engine cùng đọc một bản dữ liệu.
    \item Data Lake, vì không lớp nào sở hữu bảng; ``bảng'' chỉ là một prefix thư mục.
    \item Data Lakehouse, vì footer của file Parquet đã chứa schema.
\end{enumerate}

\item File \texttt{trips.parquet} có 4 row group. Footer cho biết khoảng giá trị của cột \texttt{fare\_amount} trong từng row group: RG1 $[2, 45]$, RG2 $[3, 60]$, RG3 $[5, 480]$, RG4 $[1, 38]$. Với truy vấn \texttt{SELECT avg(tip\_amount) FROM trips WHERE fare\_amount > 100}, một engine hỗ trợ projection pushdown và predicate pushdown sẽ đọc những gì?

\begin{enumerate}[label=\alph*), leftmargin=2.2em, itemsep=1pt]
    \item Cả 4 row group, đọc tất cả các cột.
    \item Chỉ RG3, và chỉ hai cột \texttt{fare\_amount}, \texttt{tip\_amount}.
    \item Chỉ RG3, nhưng đọc tất cả các cột.
    \item Cả 4 row group, nhưng chỉ đọc cột \texttt{tip\_amount}.
\end{enumerate}

\item Một job upstream ghi nhầm một file có cột \texttt{VendorID} kiểu chuỗi vào cùng prefix của Data Lake, nơi các file khác đều lưu \texttt{VendorID} kiểu số nguyên. Lỗi này thường được phát hiện vào lúc nào?

\begin{enumerate}[label=\alph*), leftmargin=2.2em, itemsep=1pt]
    \item Ngay lúc ghi, vì object storage kiểm tra schema của file.
    \item Ngay lúc ghi, vì Parquet không cho ghi chuỗi vào một cột số.
    \item Không bao giờ, vì engine tự bỏ qua các file sai kiểu.
    \item Khi một engine quét prefix đó, thường là trong truy vấn của một người khác.
\end{enumerate}

\item Một job ghi dữ liệu tháng 2 thành 4 file và bị dừng sau khi ghi xong file thứ hai. Một reader truy vấn ngay sau đó. Kết quả nào đúng?

\begin{enumerate}[label=\alph*), leftmargin=2.2em, itemsep=1pt]
    \item Data Lake: reader thấy 2 file đã ghi mà không có cảnh báo. Iceberg: reader vẫn thấy snapshot cũ vì con trỏ trong catalog chưa đổi.
    \item Data Lake: reader báo lỗi thiếu file. Iceberg: reader thấy 2 file đã ghi.
    \item Cả hai đều thấy 2 file đã ghi, vì object storage có strong consistency.
    \item Cả hai đều không thấy file nào, vì object storage chỉ hiển thị file khi job kết thúc.
\end{enumerate}

\item Cho năm thành phần của một bảng Iceberg: (1) \texttt{snap-<id>.avro}, (2) \texttt{\dots-m0.avro}, (3) \texttt{00006-\dots.metadata.json}, (4) \texttt{part-0.parquet}, (5) bản ghi của bảng trong catalog. Khi lập kế hoạch đọc bảng, engine đi qua chúng theo thứ tự nào?

\begin{enumerate}[label=\alph*), leftmargin=2.2em, itemsep=1pt]
    \item (3) $\rightarrow$ (5) $\rightarrow$ (1) $\rightarrow$ (2) $\rightarrow$ (4)
    \item (5) $\rightarrow$ (1) $\rightarrow$ (3) $\rightarrow$ (2) $\rightarrow$ (4)
    \item (5) $\rightarrow$ (3) $\rightarrow$ (1) $\rightarrow$ (2) $\rightarrow$ (4)
    \item (5) $\rightarrow$ (3) $\rightarrow$ (2) $\rightarrow$ (1) $\rightarrow$ (4)
\end{enumerate}

\item Writer A và writer B cùng đọc metadata phiên bản 7 của một bảng Iceberg, rồi mỗi bên ghi data file mới vào một partition khác nhau. A commit trước và thành công, con trỏ giờ trỏ tới phiên bản 8. Điều gì xảy ra khi B commit?

\begin{enumerate}[label=\alph*), leftmargin=2.2em, itemsep=1pt]
    \item B ghi đè phiên bản 8, dữ liệu của A bị mất.
    \item Phép compare-and-swap của B thất bại vì con trỏ không còn là 7. B đọc phiên bản 8, thấy không xung đột, tạo phiên bản 9 chứa cả hai thay đổi rồi commit lại.
    \item B bị chặn cho tới khi A trả lại khóa (lock) trên bảng.
    \item Catalog tự gộp hai file \texttt{metadata.json} của A và B thành một.
\end{enumerate}

\item Đoạn code PyIceberg sau chạy trên một catalog trống. Giả sử lệnh \texttt{delete} xóa ít nhất một dòng. Sau khi chạy xong, \texttt{tbl.snapshots()} có bao nhiêu snapshot?

\begin{aivncodebox}[Code]
\begin{python}
tbl = catalog.create_table("silver.trips", schema=schema)
tbl.append(jan)
with tbl.update_schema() as u:
    u.add_column("tip_per_mile", DoubleType())
with tbl.update_spec() as s:
    s.add_field("pickup_at", DayTransform(), "pickup_day")
tbl.append(feb)
tbl.delete(delete_filter=GreaterThan("fare_amount", 500.0))
\end{python}
\end{aivncodebox}
\begin{enumerate}[label=\alph*), leftmargin=2.2em, itemsep=1pt]
    \item 3
    \item 4
    \item 5
    \item 6
\end{enumerate}

\item Tiếp tục câu 7. Khi đọc các dòng tháng 1 (được ghi trước lệnh \texttt{add\_column}), cột \texttt{tip\_per\_mile} có giá trị gì?

\begin{enumerate}[label=\alph*), leftmargin=2.2em, itemsep=1pt]
    \item NULL, vì file tháng 1 không chứa field ID của cột mới và không bị ghi lại.
    \item Giá trị được Iceberg tự tính từ \texttt{tip\_amount} và \texttt{trip\_distance}.
    \item Truy vấn báo lỗi schema mismatch cho tới khi ghi lại toàn bộ bảng.
    \item 0, vì Iceberg điền giá trị mặc định của kiểu double.
\end{enumerate}

\item Một bảng Iceberg ban đầu không có partition và đang chứa 1 file dữ liệu tháng 1. Đội dữ liệu thêm partition \texttt{day(pickup\_at)} rồi append dữ liệu tháng 2 (28 ngày). Nhận định nào đúng?

\begin{enumerate}[label=\alph*), leftmargin=2.2em, itemsep=1pt]
    \item Iceberg tự ghi lại file tháng 1 thành 31 file theo ngày.
    \item Muốn được partition pruning, truy vấn phải lọc trên cột \texttt{pickup\_day}.
    \item File tháng 1 bị loại khỏi snapshot hiện hành vì không khớp partition spec mới.
    \item File tháng 1 giữ spec cũ (\texttt{spec\_id} 0), file tháng 2 theo spec mới (\texttt{spec\_id} 1); điều kiện trên \texttt{pickup\_at} vẫn loại được file ở cả hai spec.
\end{enumerate}

\item Bảng \texttt{orders} (2 TB) nhận dữ liệu CDC từ database vận hành: cứ 5 phút lại có vài nghìn lệnh \texttt{UPDATE} rải rác khắp bảng. Dashboard đọc bảng vài lần mỗi giờ. Chiến lược nào phù hợp nhất?

\begin{enumerate}[label=\alph*), leftmargin=2.2em, itemsep=1pt]
    \item Copy-on-Write, vì đọc luôn nhanh nhất.
    \item Ghi lại toàn bộ bảng sau mỗi 5 phút để không phải dùng delete file.
    \item Merge-on-Read (delete file hoặc deletion vector), kết hợp compaction định kỳ.
    \item Bỏ table format, ghi thẳng file mới vào Data Lake.
\end{enumerate}

\item Dùng PyIceberg 0.12, đội dữ liệu chạy \texttt{expire\_snapshots} chỉ giữ snapshot hiện hành, rồi thấy dung lượng thư mục \texttt{data/} không giảm. Giải thích nào đúng và cần làm gì tiếp?

\begin{enumerate}[label=\alph*), leftmargin=2.2em, itemsep=1pt]
    \item Lệnh expire chưa được commit; chạy lại lệnh là đủ.
    \item Expire chỉ xóa snapshot khỏi metadata. Các file không còn được tham chiếu đã thành orphan file và cần được dọn riêng, với ngưỡng thời gian đủ xa để không xóa nhầm file của job đang ghi.
    \item Các file vẫn cần cho time travel về những snapshot đã expire nên không được xóa.
    \item Cần chạy compaction thì dung lượng mới giảm.
\end{enumerate}

\item Dữ liệu khách hàng đến từ CRM và ERP với địa chỉ mâu thuẫn nhau. Cần khử trùng lặp, chọn bản ghi đáng tin cậy và lưu lịch sử thay đổi địa chỉ theo SCD Type 2. Trong kiến trúc Medallion, việc này nên đặt ở lớp nào?

\begin{enumerate}[label=\alph*), leftmargin=2.2em, itemsep=1pt]
    \item Silver
    \item Bronze
    \item Gold
    \item Trong chính hệ thống CRM và ERP
\end{enumerate}

\item \texttt{fact\_trips} có 5.848.609 dòng. Sau khi \texttt{INNER JOIN} với \texttt{dim\_payment\_type} (không có điều kiện \texttt{WHERE}), kết quả chỉ còn 5.712.662 dòng. Nguyên nhân khả dĩ nhất và cách phát hiện là gì?

\begin{enumerate}[label=\alph*), leftmargin=2.2em, itemsep=1pt]
    \item DuckDB làm tròn số khi JOIN; cần đổi sang kiểu \texttt{DECIMAL}.
    \item Một số \texttt{payment\_type\_key} trong fact không có trong dimension (ví dụ mã 0) nên bị \texttt{INNER JOIN} loại; phát hiện bằng test khóa ngoại mồ côi (\texttt{LEFT JOIN \dots{} WHERE dim.key IS NULL}).
    \item Snapshot Iceberg đã thay đổi giữa hai truy vấn.
    \item Dimension có dòng trùng khóa nên JOIN làm giảm số dòng.
\end{enumerate}

\item Nhận định nào \textbf{không} phải là vấn đề của kiến trúc hai tầng (Data Lake + Data Warehouse) mà Armbrust và cộng sự nêu ra?

\begin{enumerate}[label=\alph*), leftmargin=2.2em, itemsep=1pt]
    \item Dữ liệu trong warehouse trễ hơn nguồn hàng giờ đến hàng ngày.
    \item Phải lưu hai bản dữ liệu và trả tiền compute cho ETL.
    \item Thư viện ML khó đọc trực tiếp dữ liệu nằm trong warehouse.
    \item Object storage không lưu được file Parquet.
\end{enumerate}

\end{enumerate}
\end{document}
